%% Vol I, Appendix B — Word-Level ACL System
%% SOURCE: docs/working/architecture/03-architecture/word-acl/DESIGN.md
%%         .claude/CLAUDE.md (Word-Level ACL System section)
%%         proof/ACL_*.thy (5 files)

\chapter{Word-Level ACL System}
\label{vol1:app:acl}

%% TODO(bob): this appendix summarizes the word-level ACL system implemented
%% through Phase 6. Full design spec is in word-acl/DESIGN.md.

\section{Design Summary}
\label{vol1:sec:acl:design}

%% SOURCE: docs/working/architecture/03-architecture/word-acl/DESIGN.md
%% TODO(bob): brief description of the ACL model: four C fields in DictEntry,
%% two VM flags, all policy logic in ACL.4th (no new C primitives for policy).

The word-level ACL system is implemented through Phase 6. Key constraints
(cite the design doc for full details):

\begin{itemize}
  \item All policy logic in \texttt{ACL.4th} — no new C primitives for policy.
  \item Four \texttt{DictEntry} C fields: \texttt{acl\_ttl}, \texttt{acl\_allow},
        \texttt{acl\_mode}, \texttt{acl\_pinned}.
  \item Two VM flags: \texttt{emergency\_console} (fault handler active) and
        \texttt{zuse\_session} (superuser authenticated).
  \item \texttt{ACL.4th} is self-activating; \texttt{init.4th} only needs
        \texttt{S" ACL.4th" EXEC}.
  \item \texttt{ACL-PIN} is one-way; inheritance clears pin, copies mode.
\end{itemize}

\section{Phase Status}
\label{vol1:sec:acl:phases}

\begin{table}[ht]
  \centering
  \caption{ACL implementation phase status}
  \label{tab:acl-phases}
  \begin{tabular}{lll}
    \toprule
    Phase & Description & Status \\
    \midrule
    1 & C Infrastructure (\texttt{DictEntry} fields, interpreter hook, \texttt{acl\_recheck()}) & Done \\
    2 & \texttt{ACL.4th} FORTH policy words + \texttt{ACL-INIT-PRIMITIVES} + self-activation & Done \\
    3 & \texttt{capsules/zuse.4th} bootstrap superuser skeleton; CA root placeholder & Done \\
    4 & \texttt{init.4th} opt-in toggle & Done \\
    5 & POST tests (800/800) + Isabelle/HOL proofs (5 theory files) & Done \\
    6 & \texttt{EMERGENCY\_CONSOLE\_ENABLED} build flag + \texttt{vm\_fault\_handler} & Done \\
    7 & LithosAnanke parity — port to kernel context, three-arch acceptance & Pending \\
    8 & PKI / thumbdrive — Ed25519 challenge-response; user minting by zuse & Future \\
    \bottomrule
  \end{tabular}
\end{table}

\section{ACL Formal Proofs}
\label{vol1:sec:acl:proofs}

%% SOURCE: proof/ACL_Pin_Monotone.thy, ACL_Inherit_Clears_Pin.thy,
%%         ACL_TTL_Bounded.thy, ACL_Emergency_Bypass.thy, ACL_No_Escalation.thy
%% TODO(bob): one paragraph per theorem with statement and brief intuition.
%% Cite james:2025:proof:acl.

Five Isabelle/HOL theorems cover the ACL system's core safety properties:

\begin{enumerate}
  \item \textbf{Pin Monotonicity} (\texttt{ACL\_Pin\_Monotone.thy}) — once pinned,
        a word remains pinned through the system's lifetime.
        %% TODO(bob): state the theorem formally.
  \item \textbf{Inheritance Clears Pin} (\texttt{ACL\_Inherit\_Clears\_Pin.thy}) — inherited words
        carry mode but never carry pin status.
        %% TODO(bob): state the theorem formally.
  \item \textbf{TTL Bounded} (\texttt{ACL\_TTL\_Bounded.thy}) — all TTL values
        remain within defined bounds; no overflow.
        %% TODO(bob): state the theorem formally.
  \item \textbf{Emergency Bypass Correct} (\texttt{ACL\_Emergency\_Bypass.thy}) —
        the emergency console bypass does not escalate privileges.
        %% TODO(bob): state the theorem formally.
  \item \textbf{No Escalation} (\texttt{ACL\_No\_Escalation.thy}) —
        no sequence of ACL operations allows a word to acquire permissions
        it was not granted.
        %% TODO(bob): state the theorem formally.
\end{enumerate}
